\section{Data and Experimental Setup}

\subsection{Datasets}

We evaluate on two public cybersecurity datasets.

\textbf{LANL Cybersecurity Dataset (LANL-2015).}
The primary dataset is the Los Alamos National Laboratory cybersecurity dataset from 2015~\cite{lanl2015}, covering 58 days of continuous network activity across an enterprise network. It contains approximately 1.6 billion events spanning 12,425 unique users and 17,684 unique computers, organized into gzipped CSV files: \texttt{auth.txt.gz} records authentication events (timestamp, source and destination users and computers, authentication type, logon type, orientation, and success/failure status); \texttt{flows.txt.gz} records network flow events (timestamp, duration, source and destination computers and ports, protocol, packet count, and byte count); and \texttt{redteam.txt.gz} provides ground-truth labels for 749 red team attack events spanning 308 unique source-destination host pairs. The red team events include lateral movement via SMB, SSH, RDP, credential reuse, and network reconnaissance, with 93.6\% of attacks originating from a single compromised host (C17693) used as a jump box.

To scope the data to a manageable size, we apply a time-windowing strategy. For each red team event, we define a window of $\pm$3,600 seconds around the event timestamp. Overlapping windows are merged, producing 25 non-overlapping intervals that capture approximately 104 million authentication events and 31.6 million flow events.

\textbf{DAPT2020 Dataset.}
The secondary dataset is DAPT2020~\cite{dapt2020}, which provides pre-extracted CICFlowMeter features from simulated APT attack traffic. It contains 86,690 flow records with 77 numerical features including packet length statistics, inter-arrival times, flag counts, and segment size averages, along with activity and stage labels. Lateral movement samples are identified by filtering the \texttt{Stage} column for entries containing ``lateral movement.'' We use this dataset to benchmark our graph-based approach against classical anomaly detection methods that operate on tabular features.

\subsection{Evaluation Metrics}

We report four metrics, all computed in pair-space: \emph{recall} (fraction of red-team pairs detected), \emph{false positive rate} (fraction of non-red-team pairs flagged), \emph{F1 score} (harmonic mean of precision and recall), and \emph{AUC} (computed using scikit-learn's \texttt{roc\_auc\_score} over all valid edges, with labels from the red-team pair set). AUC is threshold-independent: it measures whether the scoring ranks attack edges above normal ones regardless of where you set the cutoff. Pair-space evaluation is tough here --- 308 red-team pairs among hundreds of thousands of total edges means precision will always be low, and the precision-recall tradeoff is steep.

\subsection{Baseline Methods}

We compare against two classical anomaly detection methods on DAPT2020: One-Class SVM (RBF kernel, $\gamma{=}\text{``scale''}$, $\nu{=}0.1$) and Isolation Forest (100 estimators, contamination$= 0.05$). Both are trained on normal-only samples and tested on data containing both normal and lateral movement traffic, using the 77 CICFlowMeter features. An important caveat: these baselines run on DAPT2020 while our graph methods run on LANL-2015, so the numbers are not directly comparable. Running all methods on both datasets is planned future work.

\subsection{Implementation Details}

The pipeline is implemented in Python using igraph for graph operations and scikit-learn for baselines and metrics. All hyperparameters live in a JSON config file (\texttt{pipeline\_config.json}). Several of these --- the scoring weights (Equations 1--2), the auth weight multiplier ($1.5\times$), the path boost factor ($0.1$), the hop limit (4), and the top-$k$ path count (50) --- are initial values that have not been systematically tuned. We report results with these defaults and leave proper hyperparameter optimization (grid search, Bayesian optimization) to future work. Streaming construction processes the full LANL-2015 dataset (135M+ events in the time windows) in a single pass using approximately 4 GB of RAM. End-to-end execution for all three graph variants plus baselines takes approximately 4 hours on a 12-core machine.
